% META: {"id": "1SPE-SECDEG-CO-046", "chapitre": "1SPE-SECOND-DEGRE", "type_objet": "corrige", "exercice_ref": "1SPE-SECDEG-EX-046", "capacites_codes": ["C3"], "sources_inspiration": [], "status": "generated"}
% BEGIN-VERIFY
% from sympy import symbols, solve, sqrt
% x = symbols('x')
% eq = x**2 - 40*x + 200
% Delta = (-40)**2 - 4*1*200
% assert Delta == 800
% sols = sorted(solve(eq, x), key=float)
% assert sols == [20 - 10*sqrt(2), 10*sqrt(2) + 20]
% END-VERIFY
\begin{corrige}{1SPE-SECDEG-EX-046}

\textbf{1. Mise en équation.}

\commentaireMarge{« Coût de $300$~€ » se traduit par $C(x) = 300$.}

On cherche $x$ tel que $C(x) = 300$, soit :
\[
x^2 - 40x + 500 = 300 \quad \Longleftrightarrow \quad x^2 - 40x + 200 = 0.
\]

\medskip
\textbf{2. Discriminant.}

Avec $a = 1$, $b = -40$, $c = 200$ :
\[
\Delta = b^2 - 4ac = (-40)^2 - 4 \times 1 \times 200 = 1\,600 - 800 = 800.
\]

\medskip
\textbf{3. Résolution et conclusion.}

\commentaireMarge{$\Delta > 0$ : deux solutions distinctes.}

$\Delta > 0$, donc l'équation admet deux solutions :
\[
x_1 = \frac{40 - \sqrt{800}}{2} = 20 - 10\sqrt{2} \approx 5{,}86, \qquad x_2 = \frac{40 + \sqrt{800}}{2} = 20 + 10\sqrt{2} \approx 34{,}14.
\]

Les deux solutions appartiennent à l'intervalle $[0\,;\,60]$. Dans le contexte, le nombre de bracelets étant un entier, l'entreprise atteint un coût de $300$~€ pour environ $6$ bracelets ou pour environ $34$ bracelets produits.

\end{corrige}
